\subsection{上界与下界}

定义5. 设E是一个预序集，且X是E的一个子集。任一元素 \(x \in \mathbf{E}\) ，若它使得 \(x \leqslant y\) （相应地 \(x \geqslant y\) ）对所有 \(y \in \mathbf{X}\) 成立，则被称为X在E中的下界（相应地，上界）。

X 的每个上界都是 X 关于相反序的下界，反之亦然。

如果 \(x\) 是 \(\mathbf{X}\) 的一个下界，则每个元素 \(z \leqslant x\) 也是 \(\mathbf{X}\) 的一个下界。 \(\mathbf{X}\) 的一个下界也是 \(\mathbf{X}\) 的每个子集的下界。有序集 \(\mathbf{X}\) 有最小元当且仅当存在 \(\mathbf{X}\) 的一个下界，它属于 \(\mathbf{X}\) .

预序集E的子集X的下界集合可能为空：当X=E且E是没有最小元素的有序集时，就是这种情况。

E的子集X，如果其下界（相应地，上界）集合非空，则称X下有界（相应地，上有界）。既下有界又上有界的子集称为有界。如果X下有界（相应地，上有界，有界），则X的每个子集也是如此。

由单个元素组成的每个子集都是有界的。但由两个元素组成的子集不一定上有界或下有界（第10号）。

设E是一个预序集，并设f是从任意集合A到E的一个映射。映射f（按语言的滥用）被称为下有界（相应地，上有界，有界），如果集合 \(f(\mathbf{A})\) 在E中是下有界（相应地，上有界，有界）的。

